% =============================================================
% Annexe D — Conception détaillée
% La conception couche par couche (ancienne section D.1) a été remontée
% dans le corps, en section \ref{sec:conception-couches} du chapitre 4 :
% elle n'est pas dupliquée ici. Cette annexe conserve le registre complet
% des décisions, les composants écartés et la prise en charge des vingt
% exigences.
% =============================================================
\chapter{Conception détaillée}
\label{ann:conception-detaillee}

Cette annexe développe trois ensembles que le \cref{ch:conception} ne porte qu'en
synthèse : la justification de chacune des dix-sept décisions d'architecture, le motif de rejet et
la condition de réévaluation des sept composants écartés, et l'état de validation des vingt
exigences. La conception détaillée couche par couche figure dans le corps
(\cref{sec:conception-couches}) et n'est pas dupliquée ici.

\section{Registre complet des décisions d'architecture}
\label{ann:registre-decisions}

Les dix-sept décisions D00 à D16 figurent au \cref{tab:registre-decisions-annexe} avec leur statut
--- \emph{validée} (conservée telle quelle), \emph{retouchée} (ajustée par rapport à la conception
initiale), \emph{ajoutée} ou \emph{réserve} --- et leur justification. Les deux dernières, dont le
corps porte le dossier de décision entier, y renvoient plutôt que de le répéter.

\begingroup
\small
\setlength{\tabcolsep}{4pt}
\setlength{\LTpre}{4pt}
\setlength{\LTpost}{4pt}
\begin{longtable}{|P{3.2cm}|P{12.1cm}|}
\caption[Registre complet des décisions d'architecture]{Registre complet des décisions
d'architecture, avec leur statut et leur justification}
\label{tab:registre-decisions-annexe} \\
\hline
\textbf{Décision et statut} & \textbf{Justification} \\
\hline
\endfirsthead
\hline
\textbf{Décision et statut} & \textbf{Justification} \\
\hline
\endhead
\textbf{D00} Google Cloud comme fournisseur unique du socle --- \emph{validée} &
La contrainte humaine écarte l'auto-hébergement des briques de sécurité, la contrainte économique
la facturation au forfait ; entre les trois fournisseurs restants, le critère décisif --- stockage,
langage de détection et recherche vectorielle dans un même entrepôt, ce qui supprime la base
vectorielle séparée qu'imposeraient les deux autres --- est établi au \cref{ch:etat-art}.
Décision la moins réversible du registre, elle conditionne D01 à D07 et D11 \\ \hline
\textbf{D01} Exécution de conteneurs sans serveur comme unique plateforme --- \emph{validée} &
Aucun système d'exploitation à durcir, chiffrement interne natif, identité par service, réduction à
zéro hors trafic. Un orchestrateur apporterait une souplesse inutile ici, au prix d'un cluster à
exploiter \\ \hline
\textbf{D02} Point d'entrée unique avec filtrage applicatif --- \emph{validée} &
Un point d'entrée unique est un point de contrôle unique, et le filtrage managé couvre les risques
applicatifs courants sans composant à maintenir. Corollaire assumé : l'évaluation étant
premier-match-gagne, l'ordre des règles devient une propriété de sécurité \\ \hline
\textbf{D03} Base de données managée, adresse privée, restauration à un instant donné ---
\emph{validée} &
Correspond au besoin de l'application hébergée. L'instance de recette a été basculée en haute
disponibilité régionale le 08/08/2026 (voir D13) \\ \hline
\textbf{D04} Entrepôt de données généraliste comme système de supervision --- \emph{validée} &
Le système de supervision est \emph{construit}, non acheté : c'est l'apport du projet. L'entrepôt
fournit le stockage, le langage de détection et --- point décisif --- la recherche vectorielle
nécessaire à L6 sans base spécialisée \\ \hline
\textbf{D05} Modèle de représentation non génératif pour l'enrichissement --- \emph{validée} &
Sortie plus reproductible qu'une génération libre, exécution sur processeur généraliste, surface
d'injection par instruction réduite --- sans supprimer les risques liés aux entrées, au modèle ou à
sa chaîne d'approvisionnement \\ \hline
\textbf{D06} Recherche vectorielle déportée dans l'entrepôt --- \emph{validée} &
Supprime un composant entier ; cohérent avec le principe de justification \\ \hline
\textbf{D07} Registre d'images unique, retrait du second registre --- \emph{retouchée} &
Deux registres constituent deux surfaces de chaîne d'approvisionnement et une ambiguïté sur la
source de vérité. Les images de base restent accessibles par un dépôt miroir en lecture \\ \hline
\textbf{D08} Fédération d'identité pour la chaîne de livraison --- \emph{ajoutée} &
Évite une clé de longue durée dans le chemin de livraison. La condition d'attribut porte sur le
dépôt \emph{et} sur la branche, pour la raison exposée en \cref{subsec:wif} ; la politique
d'organisation censée interdire toute création de clé est inapplicable dans le périmètre
administratif disponible (É7) \\ \hline
\textbf{D09} Détection de secrets en tête de chaîne --- \emph{ajoutée} &
Coût quasi nul, gain élevé ; l'audit du \cref{ch:cadre-existant} démontre le coût d'un
secret mal géré \\ \hline
\textbf{D10} Autorisation applicative par rôles conservée --- \emph{validée} &
L'autorisation applicative reste dans l'application, celle d'infrastructure protège
l'infrastructure : les deux plans ne se substituent pas l'un à l'autre \\ \hline
\textbf{D11} Sortie réseau par le chemin contrôlé --- \emph{réserve, levée par D15} &
Le routage était décrit en code, mais le contrôle effectif était porté par le paramètre de sortie
de la plateforme d'exécution et non par une règle de pare-feu : c'est cet écart entre configuration
établie et propriété vérifiée que le statut enregistrait. La refonte du plan réseau l'a levé ; T20
le confirme (conforme, 24/08/2026) et T14 en délimite la portée résiduelle (non conforme,
25/08/2026) \\ \hline
\textbf{D12} Infrastructure en code comme standard unique de provisionnement ---
\emph{validée, renforcée} &
Réponse directe au point bloquant B5. Trois gains : reproductibilité, contrôle avant création,
détection de dérive \\ \hline
\textbf{D13} Haute disponibilité régionale appliquée --- \emph{révision de D03, 08/08/2026} &
Bascule appliquée en recette et confirmée en direct le 08/08/2026, en 11~min~2~s : le point de
défaillance zonal de la base est fermé et le code d'infrastructure décrit cet état. Limite
subsistante : le basculement de zone n'a jamais été provoqué ni chronométré en conditions
réelles \\ \hline
\textbf{D14} Réécriture du tableau de bord sur un cadre à rendu serveur avec passerelle applicative
--- \emph{validée} &
L'interface antérieure exécutait ses requêtes depuis le poste de l'analyste et détenait pour cela
un mot de passe de service. La réécriture déplace tous les appels côté serveur : le navigateur ne
détient plus qu'un jeton de session, la clé de signature n'est jamais connue du service qui rend
les pages, et l'autorité d'autorisation reste l'API de supervision --- le secret antérieur devenant
un vestige dont le retrait accompagne la clôture de cette décision \\ \hline
\textbf{D15} Sortie réseau directe, refus par défaut en sortie et pare-feu ciblé par identité ---
\emph{appliquée du 11 au 25/08/2026} &
Le verrou levé est qu'une règle de refus n'aurait rien bloqué : le trafic ne passait pas là où elle
se serait appliquée. Plan d'adressage, douze règles, coût et limites sont établis en
\cref{subsec:reseau-cible}, la conduite de la migration en \cref{subsec:migration-reseau}.
Vérification : protocole T20, conforme le 24/08/2026 \\ \hline
\textbf{D16} Isolation des données proportionnée à la sensibilité ---
\emph{arrêtée le 29/08/2026, non appliquée} &
Justification, test d'admission d'un locataire, coût et prix payé en
\cref{subsec:isolation-donnees}. Le statut est un choix d'ordonnancement énoncé comme tel : deux
migrations d'infrastructure dans la même quinzaine ne sont pas soutenables par un opérateur
unique \\
\hline
\end{longtable}
\endgroup

Un second registre, tenu au fil de la réalisation, enregistre les quinze décisions prises du
30/07/2026 au 19/08/2026 : le premier documente le \emph{pourquoi} de l'architecture cible, le
second celui de chaque évolution.

\section{Composants évalués puis écartés}
\label{ann:composants-ecartes}

Sept éléments ont été évalués puis écartés en application du principe de justification. Les
documenter démontre que leur absence est une décision et non une lacune ; chacun figure au
\cref{tab:composants-ecartes-annexe} avec son motif de rejet et sa \emph{condition de
réévaluation}.

\begingroup
\small
\setlength{\tabcolsep}{4pt}
\setlength{\LTpre}{4pt}
\setlength{\LTpost}{4pt}
\begin{longtable}{|P{3.4cm}|P{7.0cm}|P{4.6cm}|}
\caption[Composants écartés et conditions de réévaluation]{Composants évalués puis écartés, motif et condition de réévaluation}
\label{tab:composants-ecartes-annexe} \\
\hline
\textbf{Composant} & \textbf{Motif de rejet} & \textbf{Condition de réévaluation} \\
\hline
\endfirsthead
\hline
\textbf{Composant} & \textbf{Motif de rejet} & \textbf{Condition de réévaluation} \\
\hline
\endhead
Périmètre de protection contre l'exfiltration &
Mécanisme puissant mais lourd à exploiter --- exceptions, diagnostic complexe --- et
disproportionné pour un périmètre à projet unique &
Passage en production multi-équipes avec données réglementées \\ \hline
Orchestrateur de conteneurs et maillage de services &
La plateforme retenue fournit déjà le chiffrement interne et l'identité par service ; un maillage
ajouterait un plan de contrôle entier sans couvrir de menace résiduelle identifiée &
Besoin de charges longues, à état, ou de politiques fines entre services \\ \hline
Cache distribué &
Nécessaire aux seules limitations de débit distribuées et au cache multi-instances, la mise à
l'échelle étant hors objectif (BNF7) &
Montée en charge réelle \\ \hline
Déploiement multi-région &
Coût doublé, complexité triplée, valeur de démonstration nulle ; les sauvegardes couvrent le risque
de perte de données &
Exigence contractuelle de disponibilité \\ \hline
Vérification de signature des images au déploiement &
Contrôle excellent, mais d'apport marginal après le contrôle de vulnérabilité, le registre unique
et la fédération d'identité &
Première extension recommandée après le projet \\ \hline
Passerelle d'authentification devant le tableau de bord &
L'autorisation applicative par rôles couvre le besoin du périmètre ; une passerelle ajouterait une
fédération d'entreprise redondante, et son service managé exige une organisation cloud dont le
projet ne dispose pas &
Ouverture du tableau de bord à des identités d'entreprise \\ \hline
Modèle de langue génératif pour l'analyse d'incidents &
Écarté \textbf{par conception} : sortie non déterministe et surface d'injection par le contenu des
journaux, inacceptables dans une chaîne de preuve &
\textbf{Aucune} \\
\hline
\end{longtable}
\endgroup

Le dernier rejet est de nature différente des six autres : ceux-là sont des arbitrages de coût ou
de complexité, révisables si le contexte change ; celui-ci est un contrôle de sécurité, et sa
raison ne dépend pas du contexte.

\section{Prise en charge des vingt exigences par la conception}
\label{sec:annexe-couverture-exigences}

Le \cref{tab:couverture-exigences} donne, exigence par exigence, la ou les couches qui la prennent
en charge et son état de validation à la date de référence ; le corps
(\cref{sec:couverture-exigences}) n'en retient que la synthèse et les trois réserves ---
EX7, EX9 et le couple EX16 / EX18. L'intitulé complet de chaque exigence figure dans la matrice de
traçabilité du \cref{ch:besoins-menaces}. À la date de \dateref{}, les vingt protocoles
portent tous un résultat daté ; six d'entre eux sont en outre rejoués à la cadence fixée par le
\textbf{plan de validation continue} (\cref{sec:plan-validation}).

\begingroup
\small
\setlength{\tabcolsep}{4pt}
\setlength{\LTpre}{4pt}
\setlength{\LTpost}{4pt}
\begin{longtable}{|P{1.3cm}|P{1.9cm}|P{11.8cm}|}
\caption[Prise en charge des exigences]{Prise en charge des exigences par la conception et état de validation}
\label{tab:couverture-exigences} \\
\hline
\textbf{Exigence} & \textbf{Couche(s)} & \textbf{Statut} \\
\hline
\endfirsthead
\hline
\textbf{Exigence} & \textbf{Couche(s)} & \textbf{Statut} \\
\hline
\endhead
EX1 & L1 & Prévue ; \textbf{T1 non conforme} --- motif brut traité par redirection \\ \hline
EX2 & L1, L5 & Prévue ; T2 conforme \\ \hline
EX3 & L1, L7 & Prévue ; T3 partiellement conforme \\ \hline
EX4 & L4, L5 & Prévue ; T4 partiellement conforme \\ \hline
EX5 & L2 & Prévue ; T5 conforme \\ \hline
EX6 & L2, L3 & Prévue ; T6 partiellement conforme \\ \hline
EX7 & L2 & Réserve : politique d'organisation inapplicable (écart É7) ; T7 conforme, garantie vérifiée par inventaire et non préventive \\ \hline
EX8 & F3 & Prévue ; T8 conforme --- refus réellement survenu \\ \hline
EX9 & L2, F3 & Réserve : déploiement par étiquette de commit, empreinte stricte non prouvée (écart É8) ; T9 partiellement conforme \\ \hline
EX10 & L2, F3 & Prévue ; T10 partiellement conforme \\ \hline
EX11 & L2, L5 & Prévue ; T11 conforme \\ \hline
EX12 & L7 & Prévue ; T12 partiellement conforme (22/08/2026) \\ \hline
EX13 & L7 & Prévue ; T13 partiellement conforme \\ \hline
EX14 & L4, L6 & \textbf{Ouverte} : T14 non conforme (25/08/2026), sortie managée non gouvernée \\ \hline
EX15 & L6 & Prévue ; T15 partiellement conforme (22/08/2026) \\ \hline
EX16 & F7 & Réserve : bornée aux ressources déjà gérées ; T16 partiellement conforme \\ \hline
EX17 & F7 & Prévue ; T17 conforme (20/08/2026) \\ \hline
EX18 & F7 & Réserve : revue manuelle, sans porte d'approbation outillée ; T18 partiellement conforme \\ \hline
EX19 & L7 & Prévue ; T19 partiellement conforme (21/08/2026) \\ \hline
EX20 & L4 & Journalisation des refus déployée le 19/08/2026 ; T20 conforme (24/08/2026) \\
\hline
\end{longtable}
\endgroup

\section{Conception détaillée par couche : les contrôles apportés}
\label{ann:conception-couches}

Le \cref{tab:conception-couches} reprend les sept niveaux du
\cref{tab:conception-couches-synthese}, dans le même ordre, et y ajoute la colonne que le corps ne
porte pas : les \textbf{contrôles effectivement apportés}. C'est elle, et non la liste des
composants, qui établit ce que la conception oppose à chaque frontière de confiance.

\footnotesize
\setlength{\tabcolsep}{3pt}
\renewcommand{\arraystretch}{1.1}
\begin{longtable}{|P{1.70cm}|P{4.5cm}|P{7.18cm}|P{1.6cm}|}
\caption{Conception détaillée des sept niveaux : composants, produits, contrôles et exigences}
\label{tab:conception-couches} \\
\hline
\textbf{Niveau} & \textbf{Composants et produits employés} & \textbf{Contrôles apportés} &
\textbf{Exigences} \\
\hline
\endfirsthead
\caption[]{Conception détaillée des sept niveaux (suite)} \\
\hline
\textbf{Niveau} & \textbf{Composants et produits employés} & \textbf{Contrôles apportés} &
\textbf{Exigences} \\
\hline
\endhead
\textbf{L1}\newline Périmètre d'entrée &
Répartiteur global HTTPS, trois certificats gérés, moteur de filtrage applicatif, limitation de
débit, restriction géographique. \emph{Cloud Load Balancing}, \emph{Cloud Armor} --- politique
unique attachée à tous les services publiés. La zone de noms est une dépendance externe (É10) &
HTTPS exclusivement (port 80 redirigé, TLS 1.2 au minimum). Cinq familles de règles préconfigurées
en \textbf{blocage direct}, non en observation. Restriction géographique écrite comme un
\emph{refus} dont les exemptions --- sondes, administration --- sont repliées à l'intérieur, une
autorisation prioritaire étant terminale. Limitation de débit à deux niveaux : seuil strict avec
bannissement sur les chemins d'authentification, seuil global élevé contre la saturation &
EX1, EX2, EX3 \\ \hline
\textbf{L2}\newline Identité et sécurité \emph{(plan transversal)} &
\emph{Cloud IAM} (sept identités de service), \emph{Secret Manager} (neuf secrets, dont huit
chiffrés par une clé gérée par le projet), \emph{Cloud KMS} (deux trousseaux, deux clés, rotation
de quatre-vingt-dix jours), \emph{Workload Identity Federation} &
Plan interrogé à chaque appel, quel que soit le niveau (PR2). Un secret par usage, jamais dérivé
d'un autre ; le code d'infrastructure crée les conteneurs de secrets, jamais leurs valeurs. Les deux
trousseaux sont une contrainte de plateforme, non une redondance : la réplication du coffre n'accepte
qu'une clé globale, le chiffrement de l'entrepôt exige une clé régionale. Le neuvième secret, mot de
passe du tableau de bord antérieur (D14), n'est lu par aucun code &
EX5, EX6, EX7, EX10, EX11 \\ \hline
\textbf{L3}\newline Charges de travail &
\textbf{Cinq services et trois tâches} : API de supervision, tableau de bord, service d'encodage,
interface web et interface de programmation de l'application hébergée ; enrichissement, migration
du schéma, vérification quotidienne de l'isolation des bases. \emph{Cloud Run}, \emph{Cloud
Scheduler} pour deux déclencheurs, la chaîne de livraison pour le troisième &
Entrée restreinte au répartiteur pour les quatre services publiés, adresses directes du fournisseur
neutralisées ; entrée strictement interne pour le service d'encodage, invocable par la seule
identité de la tâche d'enrichissement. Une identité de service par charge de travail. Sortie
gouvernée par le pare-feu et non plus par un paramètre de plateforme depuis la refonte du
23/08/2026. Déploiement par une étiquette égale à l'empreinte du commit, qui ne vaut pas empreinte
d'image immuable (É8) &
EX6, EX9 \\ \hline
\textbf{L4}\newline Réseau &
Un réseau virtuel en mode personnalisé, \textbf{un sous-réseau par locataire} dont un réservé au
troisième, un appairage de services privés vers la base, et \textbf{douze règles de filtrage dont
cinq autorisations}. \emph{VPC}, \emph{sortie réseau directe} et \emph{Private Service Access} &
Refus par défaut \textbf{en entrée et en sortie}, l'un et l'autre journalisés ; cinq autorisations
ciblées par identité de service et par port ; refus croisés entre locataires journalisés, qui
alimentent la chaîne de détection sans la modifier ; base joignable par une adresse privée
exclusive. Principe de la couche : \textbf{le réseau n'autorise jamais par la position, il filtre
par l'identité} --- être à l'intérieur du réseau privé ne confère aucun droit &
EX4, EX20 \\ \hline
\textbf{L5}\newline Données &
\textbf{Base applicative} (comptes, rôles, journal d'audit) et \textbf{entrepôt de supervision},
c\oe ur du dispositif de détection (\cref{fig:cycle-donnee}, \cref{fig:modele-donnees}).
\emph{Cloud SQL for PostgreSQL 15} (adresse privée seule, haute disponibilité régionale,
restauration à un instant donné sur sept jours), \emph{BigQuery} (dix tables), \emph{Cloud
Storage} (clé gérée, prévention d'accès public) &
Intégrité des preuves : l'écriture sur les journaux bruts et les détections est réservée à la
collecte et aux règles, aucune identité applicative ou d'analyse n'y accédant (EX11). Coût :
partition par jour, expiration à quatre-vingt-dix jours sur les seules tables de journaux, filtre
restrictif à la source sur trois exports sur quatre. Seuil de rattachement déclaré en tête de
requête, donc auditable et modifiable sans redéploiement (BNF6) &
EX2, EX4, EX11 \\ \hline
\textbf{L6}\newline Enrichissement sémantique &
Un service d'encodage sans état, une tâche d'enrichissement par lots, un référentiel de techniques
d'attaque vectorisé. \emph{Cloud Run} en entrée interne seule, modèle de représentation de phrases
spécialisé en cybersécurité embarqué dans l'image et exécuté hors ligne, recherche vectorielle
native de l'entrepôt --- ce qui supprime la base vectorielle séparée qu'imposerait un autre
fournisseur (D05, D06) &
Poids du modèle traités comme un artefact tiers non fiable : révision figée, format refusant
l'exécution de code à la désérialisation, empreintes consignées, mêmes contrôles de vulnérabilité
que le code. Embarqués dans l'image plutôt que téléchargés à l'exécution, ce qui retire du chemin
critique une dépendance externe et un point d'injection. Appel authentifié par jeton d'identité. Un
export quantifié a été rejeté par une porte d'acceptation \emph{codée} --- plancher de similarité
et accord de premier rang --- dont l'échec interrompt la construction de l'image &
EX14, EX15 \\ \hline
\textbf{L7}\newline Observabilité \emph{(plan transversal)} &
Quatre exports de journaux, politiques d'alerte, objectifs de service, sondes de disponibilité,
lien entre une alerte affichée et sa procédure d'incident. \emph{Cloud Logging}, \emph{Cloud
Monitoring} &
Sources collectées : journaux d'audit du fournisseur, journaux des charges de travail, verdicts du
filtrage applicatif, journaux réseau (EX13) ; trois exports sur quatre portent un filtre restrictif
à la source (EX12, BNF3). Sondes sur les points de santé (BNF1). Une procédure d'incident par
famille de technique, accessible depuis le tableau de bord : une alerte sans conduite à tenir
transfère toute la charge d'interprétation à l'analyste (BNF4). Indicateurs publiés : délai de
détection, part des détections issues du rattachement, taux d'alertes non rattachées, latence,
version du modèle (BNF6) &
EX12, EX13, EX19 \\
\hline
\end{longtable}
\normalsize
\setlength{\tabcolsep}{3pt}
\renewcommand{\arraystretch}{1.0}

\section{Plan d'adressage et règles de pare-feu du réseau en service}
\label{ann:reseau-cible}

Détail du plan réseau conçu et appliqué par la décision D15 (\cref{subsec:reseau-cible}) :
l'adressage avec la justification de chaque longueur de préfixe, puis les douze règles de pare-feu
en service. Les deux tableaux forment la référence opposable du critère d'acceptation ---
\emph{aucune règle d'autorisation ne comporte de plage source}.

{\setlength{\tabcolsep}{4pt}
\begin{longtable}{|P{2.5cm}|P{2.2cm}|P{4.0cm}|P{6.0cm}|}
\caption[Plan d'adressage du réseau du socle, en service depuis le 11/08/2026]{Plan d'adressage du
réseau du socle, en service depuis le 11/08/2026. Chaque longueur de préfixe est justifiée par le
plafond d'instances déclaré et sa marge : un préfixe posé sans raison est une dette d'adressage.}
\label{tab:adressage-cible} \\
\hline
\textbf{Sous-réseau} & \textbf{Plage} & \textbf{Ce qui l'occupe} & \textbf{Justification du préfixe} \\
\hline
\endfirsthead
\caption[]{Plan d'adressage cible (suite)} \\
\hline
\textbf{Sous-réseau} & \textbf{Plage} & \textbf{Ce qui l'occupe} & \textbf{Justification du préfixe} \\
\hline
\endhead
Locataire plateforme & \texttt{10.0.8.0/26} & Interfaces de sortie de l'API de supervision, du
tableau de bord et de la tâche d'enrichissement & 60 adresses pour un plafond déclaré de six
instances, soit dix fois la marge ; c'est aussi le préfixe minimal recommandé pour la sortie réseau
directe \\ \hline
Locataire hébergé & \texttt{10.0.8.64/26} & Interfaces de sortie de l'application hébergée et de
ses deux tâches & Même préfixe : un locataire ne doit pas hériter d'un plafond de croissance
inférieur par accident de plan \\ \hline
Troisième locataire & \texttt{10.0.8.128/26} & \textbf{Créé vide, réservé} & Un sous-réseau sans
ressource ne coûte rien : l'accueil du troisième locataire, seuil de bascule déclaré du projet,
devient une ligne de configuration et non une décision d'adressage prise dans l'urgence \\ \hline
Réserve & \texttt{10.0.8.192/26} & --- & Ferme le bloc \texttt{/24} : un \texttt{/24} vaut quatre
locataires, l'extension suivante se fait sur le \texttt{/24} voisin selon la même règle \\ \hline
Appairage de services privés & \texttt{/16} en service, \texttt{/20} pour une reconstruction &
Adresse privée de la base de données & Le \texttt{/16} en place est seize fois surdimensionné et
immuable sans recréer l'appairage ; un \texttt{/20} suffit à une instance par locataire et ne vaut
donc que pour un provisionnement neuf. Séparé du bloc des locataires : une erreur de plan ne peut
pas les faire se chevaucher \\ \hline
\multicolumn{4}{|P{15.5cm}|}{\textbf{Supprimés au cours de la migration} : la plage
\texttt{10.0.3.0/28} et le connecteur, le 15/08/2026 ; les deux sous-réseaux \texttt{10.0.1.0/24}
et \texttt{10.0.2.0/24}, qu'aucune ressource n'occupait et qu'aucun module ne consommait, le
24/08/2026 --- ce qui ferme l'écart É2.} \\
\hline
\end{longtable}}

{\setlength{\tabcolsep}{4pt}
\begin{longtable}{|P{1.0cm}|P{1.2cm}|P{4.6cm}|P{3.7cm}|P{1.5cm}|P{2.1cm}|}
\caption[Règles de pare-feu en service, ciblées par identité]{Règles de pare-feu en service depuis
le 23/08/2026. Aucune autorisation ne comporte de plage source : c'est le critère d'acceptation de
la conception.}
\label{tab:regles-cibles} \\
\hline
\textbf{Réf.} & \textbf{Sens} & \textbf{Cible --- identité de service} & \textbf{Destination ou
source} & \textbf{Ports} & \textbf{Action} \\
\hline
\endfirsthead
\caption[]{Règles de pare-feu cibles (suite)} \\
\hline
\textbf{Réf.} & \textbf{Sens} & \textbf{Cible --- identité de service} & \textbf{Destination ou
source} & \textbf{Ports} & \textbf{Action} \\
\hline
\endhead
E1--E2 & Sortie & API de supervision ; application hébergée & Adresse privée de la base &
5432 & Autoriser \\ \hline
E3 & Sortie & Tableau de bord ; tâche d'enrichissement & Adresse privée de la base & tous &
\textbf{Refuser}, journalisé \\ \hline
E4 & Sortie & API de supervision ; application hébergée ; tâche d'enrichissement & Interfaces
privées du fournisseur & 443 & Autoriser \\ \hline
E5 & Sortie & Tableau de bord & \textbf{Adresse du répartiteur, et elle seule} & 443 &
Autoriser \\ \hline
E6 & Sortie & Tâche d'enrichissement & Service d'encodage & 443 & Autoriser \\ \hline
E7--E8 & Sortie & Identités d'un locataire & Plage de l'autre locataire & tous &
\textbf{Refuser}, journalisé \\ \hline
E9 & Sortie & \emph{tout le réseau} & \texttt{0.0.0.0/0} & tous & \textbf{Refuser par défaut},
journalisé \\ \hline
I1--I2 & Entrée & Identités d'un locataire & Plage de l'autre locataire & tous &
\textbf{Refuser}, journalisé \\ \hline
I3 & Entrée & \emph{tout le réseau} & \texttt{0.0.0.0/0} & tous & \textbf{Refuser par défaut},
journalisé \\
\hline
\end{longtable}}
